\documentclass[10pt]{article}
\usepackage{lingmacros}
\usepackage{tree-dvips}
\usepackage{makecell}
\begin{document}
\pagenumbering{gobble}
\section*{Misura}
La pressione è una grandezza scalare perché prendiamo esclusivamente la componente perpendicolare alla superficie della forza (vettoriale) con cui viene calcolata. \\
Femto: $10^{-12}$
\section*{Newton}
Newton è il primo a capire che non è necessaria una causa per un moto costante, ma solo per un moto di accelerazione.\\
\textbf{I principio} (inerzia) definisce i sistemi di riferimento inerziali, se nessuna forza agisce su un corpo questo si trova in quiete o moto costante.\\
\textbf{II principio} (F=$m*a$) L'accelerazione è direttamente proporzionale alla forza agente, ha validità solo nei sistemi di riferimento inerziali (non è inerziale un movimento di inerzia dovuto ad un'accelerazione).\\
\textbf{III principio} (azione-reazione) Se un corpo agisce su un altro con una forza, allora anche l'altro agirà sul primo con una forza uguale ma opposta (non si annullano perché sono su corpi diversi)(terra-luna).
\section*{Moto circolare}
L'accelerazione in una curva aumenta con il restringimento del raggio per la formula a=$\frac{v^2}{r}$; la forza deriva, nel caso di una macchina, dall'attrito tra pneumatici e asfalto. \\
Con un sistema di molle si può verificare l'inerzia e l'accelerazione di un corpo in tutti i versi e direzioni, nel nostro corpo questo viene gestito dai sacculi e i canali semicircolari del \textbf{sistema vestibolare}. Le cellule ciliate contenute nelle cupole dei canali semicircolari percepiscono l'inerzia del fluido che scorre nei canali e invia l'informazione al cervello. In base alla direzione in cui vengono piegate le ciglia si aprono più o meno i canali K$^+$ aumentando o diminuendo la quantità di neurotrasmettitori.\\
Nelle saccule gli otoliti (sali di Ca) si muovo in base alla gravità e all'accelerazione nelle macule.
\section*{Gravità}
Forza generata dall'attrazione di due corpi l'uno all'altro in base alla loro massa gravitazione, dipende dalla distanza dei corpi.
\section*{Elasticità}
Nelle \textbf{fibre collagene} la costante k aumenta con la deformazione, opponendosi e irrigidendosi nel caso di forze contrarie, mentre in una molla normale k rimane uguale e per aumentare la deformazione basta aumentare la forza applicata. \\
Nel caso di corpi estesi conta più la pressione applicata che la semplice forza, quindi P=$\frac{F}{A}$\\
La deformazione viene valutata rispetto alla forma iniziale, la differenza tra forma iniziale e successiva alla deformazione si chiama Strain. Questo valore permette di calcolare il modulo di Young con la formula $\frac{Stress (\frac{F}{A})}{Strain(\frac{\Delta L}{L})}$. Nel nostro corpo l'osso ha un modulo di Young più alto rispetto a organi come il cervello.
\section*{Leve e equilibrio}
L'equilibrio in una leva si raggiunge quando $F_1*b_1=F_2*b_2$, per questo nel caso $F_1$ dovesse essere molto grande allora si può equilibrare con una forza minore ad una distanza molto maggiore dal fulcro.\\
Se $bm>br$ sono vantaggiose perché richiedono meno forza motrice. Nel caso del nostro \textbf{braccio} siamo in una situazione di costante svantaggio quando vogliamo sollevare qualcosa perché i tendini sono in prossimità del fulcro (gomito), questo avviene perché essendo più vicino al fulcro, un movimento minimo del muscolo provoca una rotazione maggiore. \\
Carrucole e paranchi???
\section*{Lavoro}
Trasferimento di energia, tenendo sospeso un oggetto per la fisica il lavoro è pari a 0 perché è assente movimento, ma a livello biologico si consuma ATP per actina e miosina. 
\end{document}